\section{Verification and Numerical Experiments}
\label{sec:results}

This section interprets the saved measurements from the pinned Project 2
repository. The core method is the conservative FDM; the conforming
$\mathcal P_1$ FEM supplies the required depth comparison. All reported
errors use the exact benchmark solutions specified in the problem pack
\cite{problempack}. The input data, assembly, solver, and post-processing
definitions are given in Sections~2--3.

\subsection{Checks before interpreting error tables}
The smooth forcing was derived analytically from
$u_*=\sin(\pi x)\sin(\pi y)$ and $a=1+x$. An independent derivative check in
the validation suite found a maximum absolute source defect of
$7.11\times10^{-15}$ in its complex-step comparison. An affine patch with
asymmetric nonzero values on all four sides gave a maximum difference of
$4.89\times10^{-15}$ between the production boundary elimination and an
independently assembled dense reference. These checks probe the source sign,
index ordering, and Dirichlet right-hand side rather than merely reproducing
the same code path.

For the smooth $N=32$ FDM and FEM systems, the tolerance sweep in
Table~\ref{tab:alg_tol} uses the same assembled matrix and right-hand side at
each solver tolerance. Tightening the requested CG relative residual from
$10^{-6}$ to $10^{-12}$ changed the CG--LU maximum solution difference from
$5.81\times10^{-10}$ to $5.33\times10^{-15}$ for FDM and from
$4.00\times10^{-8}$ to $1.01\times10^{-14}$ for FEM. Both values at the
tightest setting lie far below the approximately $8\times10^{-4}$ nodal
discretisation error. Across all four smooth levels, the recorded
algebraic-to-discretisation error ratio remained below $10^{-3}$.

The assembled FDM matrices have zero measured relative symmetry defect in
the saved runs, positive diagonals, non-positive off-diagonals, and positive
smallest eigenvalues in the independently checked small and medium systems.
The energy identity in Theorem~\ref{thm:spd_mmatrix} is the proof of positive
definiteness; an eigensolver estimate is corroboration. The layered solutions
stay within the prescribed boundary range $[0,1]$. This supports the
discrete maximum-principle prediction, although it says nothing about the
accuracy of the interface resistance.

\subsection{Smooth coefficient: convergence and method comparison}
Table~\ref{tab:alg_smooth} reports four refinements for both discretisations,
with $N=16,32,64,128$ and $225$ to $16\,129$ unknowns. At the finest level,
FDM gives $E_{2,h}=2.491\times10^{-5}$ and
$E_\infty=5.007\times10^{-5}$; FEM gives
$2.527\times10^{-5}$ and $5.053\times10^{-5}$ in the same nodal norms.
The independent validation fit over $N=32,64,128$ is $2.00034$ for FDM and
$1.99964$ for FEM in $E_{2,h}$. For the common reconstructed continuous
$L^2$ norm, the corresponding fits are $1.99996$ and $1.99898$. The observed
orders therefore match the second-order prediction on this smooth problem.
The $a=1+x$ coefficient is affine, so arithmetic averaging equals exact
midpoint sampling; this benchmark does not test the distinct arithmetic
averaging error for a curved smooth coefficient.

The matched comparison uses the same square mesh, nodal data, number of
interior unknowns, CG stopping rule, and common reconstruction for continuous
$L^2$. At $N=128$, the latter error is $4.360\times10^{-5}$ for FDM and
$8.415\times10^{-5}$ for FEM. The two methods are both second order, but
their constants differ under this particular reconstruction and triangular
quadrature. The FEM's mathematical $L^2$ convergence rate and this measured
constant should not be conflated with its nodal error, which is much closer
to FDM's on this example.

The cost table in Section~3 reports local medians over three serial runs.
At $N=128$, FDM assembly took $2.74\times10^{-3}$ seconds and FEM assembly
$1.03\times10^{-1}$ seconds; the respective Jacobi-CG solves took
$1.53\times10^{-2}$ and $3.47\times10^{-2}$ seconds, with $192$ and $443$
iterations. The sparse LU solves were about $3.06\times10^{-2}$ and
$3.08\times10^{-2}$ seconds, respectively. Both matrices had $80\,137$
stored nonzeros after explicit zero removal. These timings are measurements
for the supplied implementations and hardware, not general speed constants
for FDM or FEM. The implemented FEM assembly performs element quadrature and
therefore has more work per unknown than this vectorised FDM assembly.

\subsection{Layered interface: accuracy, flux, and maximum principle}
For $a_L=1$ and $a_R=10$, the exact flux is $J_*=-20/11$. Odd FDM meshes
place the interface halfway between columns; even FEM meshes align it with
element edges. Table~\ref{tab:interface_refinement} extends the single-mesh
comparison of Table~\ref{tab:alg_interface} to all four odd FDM meshes.

\begin{table}[H]
\centering\small
\begin{tabular}{rrrrrr}
\toprule
$N$ & Rule & $E_{2,h}$ & $E_\infty$ & $E_q^{\max}$ & $E_q^{\mathrm{mid}}$\\
\midrule
31 & arithmetic & 6.767e-03 & 1.938e-02 & 6.629e-01 & 4.697e-02\\
31 & harmonic & 4.940e-16 & 1.443e-15 & 2.953e-14 & 1.088e-14\\
63 & arithmetic & 3.271e-03 & 9.600e-03 & 6.619e-01 & 2.294e-02\\
63 & harmonic & 1.747e-15 & 7.327e-15 & 1.081e-13 & 1.081e-13\\
127 & arithmetic & 1.607e-03 & 4.777e-03 & 6.616e-01 & 1.134e-02\\
127 & harmonic & 6.261e-15 & 1.399e-14 & 3.813e-13 & 7.372e-14\\
255 & arithmetic & 7.963e-04 & 2.383e-03 & 6.615e-01 & 5.641e-03\\
255 & harmonic & 1.085e-13 & 2.183e-13 & 4.696e-13 & 5.795e-14\\
\bottomrule
\end{tabular}

\caption{Independent layered FDM validation at contrast $c=10$.
$E_q^{\max}$ is the maximum error of the shared interface-face flux over the
interior face samples; $E_q^{\mathrm{mid}}$ is the error at the sample nearest
$y=1/2$. The harmonic method's nodal errors are at round-off.}
\label{tab:interface_refinement}
\end{table}

Harmonic FDM reproduces the nodal exact solution to a maximum error below
$2.2\times10^{-13}$ at all listed meshes. On $N=63$ the interface flux error
is $1.08\times10^{-13}$. The fitted FEM at $N=64$ likewise has a nodal
maximum error of $3.31\times10^{-14}$ and interface flux error
$5.17\times10^{-13}$. The P1 result is exact in this benchmark because the
piecewise-linear exact solution belongs to the fitted finite-element space;
general P1 solutions need not have pointwise continuous element flux.

Arithmetic FDM gives $E_\infty=9.600\times10^{-3}$ at $N=63$, falling to
$2.383\times10^{-3}$ at $N=255$. Its successive measured nodal orders are
$0.991$, $0.995$, and $0.998$, consistent with first-order convergence in
this test. The maximum shared interface-flux error, however, is approximately
$0.662$ on each mesh. The midline error decreases from $4.70\times10^{-2}$
at $N=31$ to $5.64\times10^{-3}$ at $N=255$. The maximum persists because
the top and bottom boundaries retain the prescribed exact trace while the
arithmetic interface resistance perturbs nearby interior nodes. Thus the
one-dimensional resistance calculation in Section~2 predicts the wrong
conductance but does not describe the maximum face-flux error on the
two-dimensional Dirichlet square.

The arithmetic matrix is still symmetric and conserves one shared face flux
between adjacent cells. At $N=63$, the two-sided reconstructed flux
difference at the nearest pure-material edges is only
$3.85\times10^{-13}$, while the interface flux is wrong by $0.662$ at its
worst sample. This illustrates why a small reconstructed jump or a matrix
symmetry check cannot by itself validate physical interface accuracy. Both
arithmetic and harmonic solutions remain in $[0,1]$, as predicted by their
common positive-face M-matrix structure.

\subsection{Mesh size, contrast, and iterative work}
Table~\ref{tab:conditioning_corrected} uses the independent spectral
validation record. The original condition-number diagnostic used a fixed
structured Lanczos start that missed the largest eigenvalue when $c=1$;
the independent calculation used an alternative start and checked the
constant-coefficient formula. The $c=10$ and $c=100$ entries agreed with
the repository's estimates to reported precision.

\begin{table}[H]
\centering\small
\begin{tabular}{rrrrr}
\toprule
$N$ & $c$ & $\kappa_2(A_h)$ & CG iterations & Jacobi--CG iterations\\
\midrule
15 & 1 & 9.052e+01 & 48 & 48\\
31 & 1 & 3.888e+02 & 106 & 106\\
63 & 1 & 1.608e+03 & 216 & 216\\
15 & 10 & 3.996e+02 & 99 & 52\\
31 & 10 & 1.726e+03 & 250 & 112\\
63 & 10 & 7.149e+03 & 521 & 224\\
15 & 100 & 3.628e+03 & 188 & 52\\
31 & 100 & 1.567e+04 & 576 & 112\\
63 & 100 & 6.490e+04 & 1354 & 226\\
\bottomrule
\end{tabular}

\caption{Condition number and CG iterations for harmonic FDM under mesh
refinement and contrast. Residual tolerance is $10^{-12}$; the spectral
estimates use the independently checked largest eigenvalue.}
\label{tab:conditioning_corrected}
\end{table}

At $N=63$, raising contrast from $1$ to $10$ to $100$ raises the measured
condition number from about $1.61\times10^3$ to $7.15\times10^3$ to
$6.49\times10^4$. Unpreconditioned CG iterations rise from $216$ to $521$
to $1354$. With Jacobi scaling the corresponding counts are $216$, $224$,
and $226$. At fixed contrast, refinement also increases the condition
number, consistent with the $O(h^{-2})$ upper bound. These results support
the solver choice and show the practical value of preconditioning for this
layered geometry; they do not claim contrast-independent performance for
arbitrary heterogeneous media.

